\subsection{Iteradores y funciones integradas}
\label{alg:2-1}

Estas funciones permiten recorrer estructuras de datos, transformar secuencias y escribir código más corto y legible sin sacrificar rendimiento.

\subsubsection{enumerate()}

Recorrer un iterable obteniendo simultáneamente el índice y el elemento, evitando mantener un contador manual.

\begin{lstlisting}[language=Python]
arr = [5, 8, 2]
for i, x in enumerate(arr):
    print(i, x)
for i, x in enumerate(arr, 1):     # indices comienzan en 1
    print(i, x)
\end{lstlisting}

Es especialmente útil cuando el problema utiliza índices 1-indexados.

\vspace{0.3cm}

\subsubsection{zip()}

Recorrer múltiples iterables en paralelo. El recorrido termina cuando el iterable más corto finaliza.

\begin{lstlisting}[language=Python]
a = [1, 2, 3]
b = [4, 5, 6]
for x, y in zip(a, b):
    print(x, y)

pares = list(zip(a, b))
# [(1,4), (2,5), (3,6)]
\end{lstlisting}

También puede utilizarse para desempaquetar listas de tuplas.

\begin{lstlisting}[language=Python]
pares = [(1, 4), (2, 5), (3, 6)]
a, b = zip(*pares)
print(a); print(b)
\end{lstlisting}

\vspace{0.3cm}

\subsubsection{iter() y next()}

iter() crea un iterador sobre cualquier iterable.
next() consume el siguiente elemento del iterador.
Se suelen usar cuando se procesan tokens manualmente.

\begin{lstlisting}[language=Python]
data = ["3", "10", "20", "30"]          # n y luego n valores
it = iter(data)
n = int(next(it))
arr = [int(next(it)) for _ in range(n)] # [10, 20, 30]
\end{lstlisting}

También puede especificarse un valor por defecto.

\begin{lstlisting}[language=Python]
x = next(it, None)
\end{lstlisting}

\vspace{0.3cm}

\subsubsection{any() y all()}

Evaluar rápidamente condiciones sobre un iterable.

any() devuelve True si al menos un elemento cumple la condición.
all() devuelve True únicamente si todos los elementos la cumplen.
Ambas realizan \textit{short-circuit}, por lo que dejan de recorrer el iterable tan pronto conocen la respuesta.

\begin{lstlisting}[language=Python]
if any(x % 2 == 0 for x in arr):
    print("Existe un numero par")
if all(x > 0 for x in arr):
    print("Todos son positivos")
\end{lstlisting}

\vspace{0.3cm}

\subsubsection{map()}

Aplica una función a todos los elementos de un iterable.

Es muy utilizado durante la lectura de entrada.

\begin{lstlisting}[language=Python]
nums = list(map(int, input().split()))
cuadrados = list(map(lambda x: x * x, nums))
\end{lstlisting}

También puede utilizarse con funciones propias.

\begin{lstlisting}[language=Python]
def f(x):
    return x + 1
nuevo = list(map(f, nums))
\end{lstlisting}

\vspace{0.3cm}

\subsubsection{sorted()}

Devuelve una nueva lista ordenada sin modificar la original.

Puede ordenarse mediante una función key.

\begin{lstlisting}[language=Python]
nums = sorted(nums)
nums = sorted(nums, reverse=True)
personas = sorted(personas, key=lambda x: x.edad)
pares = sorted(freq.items(), key=lambda kv: -kv[1])
\end{lstlisting}

\graybold{Complejidad:} \bigO{n \log n}

Python utiliza Timsort, un algoritmo estable muy eficiente para datos parcialmente ordenados.

\vspace{0.3cm}

\subsubsection{reversed()}

Permite recorrer una secuencia en orden inverso sin crear una copia.

\begin{lstlisting}[language=Python]
for x in reversed(arr):
    print(x)
\end{lstlisting}

Es preferible a lo siguiente, que crea una copia de la lista:

\begin{lstlisting}[language=Python]
arr[::-1]
\end{lstlisting}

\vspace{0.3cm}

\subsubsection{encode() y decode()}

encode() convierte un string en bytes; decode() hace lo contrario. En este libro
aparecen todo el tiempo, porque la lectura rápida (sección 1) usa
\texttt{sys.stdin.buffer}, que entrega \texttt{bytes} y no \texttt{str}.

\begin{lstlisting}[language=Python]
data = sys.stdin.buffer.read().split()   # lista de bytes: [b'3', b'abc', ...]
n = int(data[0])                         # int() acepta bytes directamente
s = data[1].decode()                     # decode() solo si se necesitan metodos de str
if data[2] == b"L": ...                  # comparar bytes con bytes, sin decodificar
\end{lstlisting}

Dos detalles que causan errores silenciosos: un \texttt{bytes} nunca es igual a un
\texttt{str} (\texttt{b"L" == "L"} es \texttt{False}), e indexar un \texttt{bytes}
devuelve un entero, no un carácter (\texttt{b"abc"[0]} es \texttt{97}). Por eso las
entradas del libro comparan contra códigos ASCII o contra \texttt{b"..."}.

\vspace{0.3cm}

\subsubsection*{Resumen}

\begin{center}
\begin{tabular}{|l|l|}
\hline
Función & Uso principal \\ \hline
enumerate() & Recorrer con índice \\ \hline
zip() & Recorrer varios iterables \\ \hline
iter() & Crear un iterador \\ \hline
next() & Consumir siguiente elemento \\ \hline
any() & Existe algún elemento que cumpla \\ \hline
all() & Todos cumplen \\ \hline
map() & Aplicar función a todos los elementos \\ \hline
sorted() & Ordenar sin modificar el original \\ \hline
reversed() & Recorrer en orden inverso \\ \hline
encode() & String $\rightarrow$ bytes \\ \hline
decode() & Bytes $\rightarrow$ string \\ \hline
\end{tabular}
\end{center}
